\documentclass[a4paper, UTF-8]{article}
\usepackage{algorithm}
\usepackage{algpseudocode}
\usepackage{amsmath}
\usepackage{amsfonts}
\usepackage{amsthm}
\usepackage{tikz}
\usetikzlibrary{automata, positioning, arrows}
\usepackage[utf8]{inputenc} % allow utf-8 input
\usepackage[T1]{fontenc}    % use 8-bit T1 fonts
\usepackage{hyperref}       % hyperlinks
\usepackage{url}            % simple URL typesetting
\usepackage{booktabs}       % professional-quality tables
\usepackage{amsfonts}       % blackboard math symbols
\usepackage{nicefrac}       % compact symbols for 1/2, etc.
\usepackage{microtype}      % microtypography
\usepackage{xcolor}         % colors

\theoremstyle{plain}
\newtheorem{theorem}{\indent Theorem}[section]
\newtheorem{lemma}[theorem]{\indent Lemma}
\newtheorem{assumption}{\indent Assumption}[section]
\newtheorem{note}[theorem]{\indent Notation}
\newtheorem{proposition}[theorem]{\indent Proposition}
\newtheorem{corollary}[theorem]{\indent Corollary}
\newtheorem{definition}{\indent Definition}[section]
\newtheorem{example}{\indent Example}[section]
\newtheorem{remark}{\indent Remark}[section]
\newenvironment{solution}{\begin{proof}[\indent\bf Solution]}{\end{proof}}
\renewcommand{\proofname}{\indent\bf Proof}

\begin{document}

\title{\bf{Lab4 Answers}}

\author{Yanqiao Chen*\\ 
  Department of Computer Science and Engineering\\
  Southern University of Science and Technology\\
  \texttt{12412115@mail.sustech.edu.cn}
  }

\maketitle

\section{Question 1} 

\subsection{a)} 

\begin{proof}
Let the pumping length be $p$.
Considering $w = 1^p 2^p 3^p$, by the pumping lemma we can divide it into $w = xyz$ where $|xy| \leq p$ and $|y| = l > 0$.
Clearly, we can only have $x = 1^k$ and $y = 1^l$, where $k + l \leq p$.
Thus, $z = 1^{p-k-l} 2^p 3^p$.
Now we can pump $y$ to get $x y^2 z = 1^k 1^{2l} 1^{p-k-l} 2^p 3^p = 1^{p + l} 2^p 3^p$.
Since $l > 0$, the number of 1s strictly exceeds the number of 2s and 3s. However, this string is not in the language, which is a contradiction. Therefore, the language is not regular.
\end{proof} 

\subsection{b)}

\begin{proof} 
Let the pumping length be $p$. Considering $w = (a^{p}b^p)^3$, its length is $6p$.
By the pumping lemma, we divide it into $xyz$ with $|xy| \leq p$ and $|y| = l > 0$.
This means $y$ must consist entirely of $a$'s from the first block, so $y = a^l$.
Then $xy^2z$ will be $a^{p + l} b^p a^p b^p a^p b^p$.
Since $l > 0$, the first block of $a$'s is longer than the other two blocks of $a$'s. This string cannot be divided into three identical substrings $w'w'w'$.
Therefore, the language is not regular.
\end{proof}

\subsection{c)} 

\begin{proof} 
Let the pumping length be $p \geq 2$.
Considering $w = a^{2^p}$, we divide it into $w = xyz$ where $|xy| \leq p$ and $|y| = l > 0$.
Since $w$ only consists of $a$'s, pumping $y$ simply increases the length of the string by $l$. Let's pump up by setting $i=2$.
The length of the new string $xy^2z$ is $2^p + l$.
Since $1 \leq l \leq |xy| \leq p$, we have:
$$2^p < 2^p + l \leq 2^p + p < 2^{p+1}$$
The last inequality holds because $p < 2^p$ for all $p \geq 2$. 
Thus, the length of $xy^2z$ is strictly between two consecutive powers of 2, which means $xy^2z$ is not in the language. This is a contradiction. Hence the language is not regular.
\end{proof}

\section{Question 2} 

In Example 1, we say that $0^p 1^p$ cannot be pumped since $xy^2 z \notin \{0^k1^k|k \geq 0\}$.
However, $xy^iz \in 0^*1^*$, and to prove that $0^*1^*$ is not regular, you should prove that $xy^iz$ is not in $0^*1^*$ but not $xy^2z \notin \{0^k1^k|k \geq 0\}$.
The argument cannot be directly applied to $0^*1^*$ since $0^p 1^p$ is in $0^*1^*$ and pumping its $0$s still results in a string belonging to $0^*1^*$.

\section{Question 3} 

\subsection{a)} 

\begin{proof} 
Suppose that the language is regular. Let the pumping length be $p$.
Consider $w = 0^p1^p0^p$. By the pumping lemma, we divide it into $xyz$ where $|xy| \leq p$ and $|y| = l > 0$.
This implies $x$ and $y$ consist entirely of the leading $0$s. Let $y = 0^l$.
Then $xy^2z = 0^{p + l}1^p0^p$. Since $l > 0$, the number of leading $0$s is strictly greater than the number of trailing $0$s.
This string is not in the language, which is a contradiction. Therefore, the language is not regular.
\end{proof} 

\subsection{b)} 

\begin{proof} 
Suppose for contradiction that $L = \{0^m1^n | m\ne n \}$ is a regular language.
We know that the language $L_{01} = \{0^k1^j | k, j \geq 0\}$, which is generated by $0^*1^*$, is regular.
Since regular languages are closed under complement and intersection, if $L$ is regular, then $\overline{L} \cap L_{01}$ must also be regular.
However, $\overline{L} \cap L_{01} = \{0^m1^n | m = n\}$.
We already know that $\{0^m1^n | m = n\} = \{0^k1^k | k\geq 0\}$ is not regular, which leads to a contradiction.
Therefore, $\{0^m1^n | m\ne n \}$ is not regular.
\end{proof}

\section{Question 4} 

\begin{proof}
We claim that the number of occurrences of ``01'' equals the number of occurrences of ``10'' in a string $w$ if and only if either $w \in \{\varepsilon, 0, 1\}$ or the first and last characters of $w$ are the same.
To see this, note that each occurrence of ``01'' corresponds to a transition from 0 to 1, and each ``10'' corresponds to a transition from 1 to 0. These transitions must alternate.
If $w$ starts and ends with the same character, the total number of transitions is even, so the number of ``01''s equals the number of ``10''s.
If $w$ starts and ends with different characters, the number of transitions is odd, so the counts differ by 1.

Therefore, the language can be expressed by the following regular expression:
$$D = \{\varepsilon, 0, 1\} \cup 0\Sigma^*0 \cup 1\Sigma^*1$$
Since $D$ can be described by a regular expression, it is a regular language.

\begin{remark}
    For example, $01100110$ is in the language.
    We can observe that any sequence of 0s starts with a substring $10$ and ends with a substring $01$, and any sequence of 1s starts with a substring $01$ and ends with a substring $10$.
    Thus, the number of occurrences of $01$ equals the number of occurrences of $10$.
\end{remark}

\end{proof}

\section{Question 5} 

\subsection{a)} 
\begin{proof}
We claim that $B = \{1w \mid w \in \{0,1\}^* \text{ and } w \text{ contains at least one } 1\}$.

Proof of claim:
\begin{itemize}
    \item ($\subseteq$) If $u \in B$, then $u = 1^k y$ where $y$ has at least $k \geq 1$ ones. Then $u = 1 \cdot (1^{k-1}y)$ and $1^{k-1}y$ contains at least $(k-1)+k \geq 1$ ones.
    \item ($\supseteq$) If $u = 1w$ where $w$ has at least one 1, then $u = 1^1 \cdot w$ and $w$ has at least $1$ one, so $u \in B$.
\end{itemize}

Therefore $B = 1(0\cup 1)^*1(0\cup 1)^*$, which is regular.
\end{proof}

\subsection{b)} 

\begin{proof}
We prove it by the pumping lemma. Suppose $C$ is regular. Let the pumping length be $p$.

Consider $w = 1^p 0 1^p$. Clearly, $w \in C$ because it can be written as $1^k y$ with $k=p$ and $y=01^p$, where $y$ contains exactly $p$ ones (which is at most $k$).

By the Pumping Lemma, we can divide $w = xyz$ such that $|xy| \leq p$ and $|y| = l > 0$.

Since the first $p$ characters are all $1$'s, $y$ must consist entirely of $1$'s, so $y = 1^l$.
Now, consider pumping down by setting $i=0$. The resulting string is $xz = 1^{p-l} 0 1^p$.
For $xz$ to be in $C$, it must be of the form $1^k y'$ where $y'$ contains at most $k$ ones.
The maximum number of consecutive $1$'s at the beginning of $xz$ is $p-l$. Thus, any valid prefix $1^k$ must have $k \leq p-l$.
The remaining suffix $y'$ will inevitably contain the $p$ ones that appear after the $0$.
Since $l \geq 1$, we have $p > p-l \geq k$. 

Therefore, the number of ones in $y'$ is strictly greater than $k$, meaning $xz \notin C$.
This is a contradiction. Hence the language $C$ is not regular.  
\end{proof}
 
\end{document}